% Auto-generated by EcoCivic Academic Monte Carlo Simulation Engine (N=10,000, T=12 Months)
% Target Venues: IEEE Access / MDPI Sustainability (Q1/Q2)
\begin{table*}[htbp]
\centering
\caption{Comparative Performance and Cost Analysis Across 10,000 Households Over a 12-Month Period}
\label{tab:ecocivic_comparative_analysis}
\begin{tabular}{lrrrrrr}
\toprule
\textbf{Architecture / Paradigm} & \textbf{CAPEX (\$)} & \textbf{12M OPEX (\$)} & \textbf{Total Cost (\$)} & \textbf{Fraud Recall (\%)} & \textbf{Fraud Leakage (\$)} & \textbf{Cost Reduction (\%)} \\
\midrule
Scenario 1: Full-IoT (AMI + RFID) & \$1,850,000 & \$288,000 & \$2,138,000 & 94.5\% & \$12,500 & Baseline (IoT) \\
Scenario 2: Manual Municipal Audit & \$15,000 & \$556,800 & \$571,800 & 31.8\% & \$148,200 & Baseline (Manual) \\
\textbf{Scenario 3: Proposed EcoCivic Hybrid} & \textbf{\$15,000} & \textbf{\$123,948} & \textbf{\$138,948} & \textbf{82.5\%} & \textbf{\$7,680} & \textbf{-93.5\% vs IoT / -75.7\% vs Man.} \\
\bottomrule
\end{tabular}
\vspace{1mm}
\\ \footnotesize{\textit{Note:} In Scenario 3, an additional \$143,586.00 in slashed USDC collateral was recovered from malicious actors, reducing the net municipal cost to \$0.00. Citizen fraud attempts converged from 13.59\% in Month 1 down to 0.34\% in Month 12 (-97.5\% deterrence convergence).}
\end{table*}
